\begin{register}{H}{GPIO\_OUT\_REG}{0x0004}\label{regdesc:GPIOOUTREG}
\regfield{}{32}{0}{xxxxxxxxxxxxxxxxxxxxxxxxxxxxxxxx}%
\reglabel{Reset}\regnewline%
\begin{regdesc}\begin{reglist}
\item [GPIO\_OUT\_REG]GPIO0-31 输出值。（读/写）
\end{reglist}\end{regdesc}
\end{register}


\begin{register}{H}{GPIO\_OUT\_W1TS\_REG}{0x0008}\label{regdesc:GPIOOUTW1TSREG}
\regfield{}{32}{0}{xxxxxxxxxxxxxxxxxxxxxxxxxxxxxxxx}%
\reglabel{Reset}\regnewline%
\begin{regdesc}\begin{reglist}
\item [GPIO\_OUT\_W1TS\_REG]GPIO0-31 输出置位寄存器。每一位置 1，则 GPIO\_OUT\_REG 中的相应位也会置 1。（只写）
\end{reglist}\end{regdesc}
\end{register}


\begin{register}{H}{GPIO\_OUT\_W1TC\_REG}{0x000c}\label{regdesc:GPIOOUTW1TCREG}
\regfield{}{32}{0}{xxxxxxxxxxxxxxxxxxxxxxxxxxxxxxxx}%
\reglabel{Reset}\regnewline%
\begin{regdesc}\begin{reglist}
\item [GPIO\_OUT\_W1TC\_REG]GPIO0-31 输出清零寄存器。每一位置 1，则 GPIO\_OUT\_REG 中的相应位会清零。（只写）
\end{reglist}\end{regdesc}
\end{register}


\begin{register}{H}{GPIO\_OUT1\_REG}{0x0010}\label{regdesc:GPIOOUT1REG}
\regfield{(reserved)}{24}{8}{000000000000000000000000}%
\regfield{GPIO\_OUT\_DATA}{8}{0}{xxxxxxxx}%
\reglabel{Reset}\regnewline%
\begin{regdesc}\begin{reglist}
\item [GPIO\_OUT\_DATA]GPIO32-39 输出值。（只读）
\end{reglist}\end{regdesc}
\end{register}


\begin{register}{H}{GPIO\_OUT1\_W1TS\_REG}{0x0014}\label{regdesc:GPIOOUT1W1TSREG}
\regfield{(reserved)}{24}{8}{000000000000000000000000}%
\regfield{GPIO\_OUT\_DATA}{8}{0}{xxxxxxxx}%
\reglabel{Reset}\regnewline%
\begin{regdesc}\begin{reglist}
\item [GPIO\_OUT\_DATA]GPIO32-39 输出值置位寄存器。每一位置 1，则 GPIO\_OUT1\_DATA 中的相应位也会置 1。（只读）
\end{reglist}\end{regdesc}
\end{register}


\begin{register}{H}{GPIO\_OUT1\_W1TC\_REG}{0x0018}\label{regdesc:GPIOOUT1W1TCREG}
\regfield{(reserved)}{24}{8}{000000000000000000000000}%
\regfield{GPIO\_OUT\_DATA}{8}{0}{xxxxxxxx}%
\reglabel{Reset}\regnewline%
\begin{regdesc}\begin{reglist}
\item [GPIO\_OUT\_DATA]GPIO32-39 输出值清零寄存器。每一位置 1，则 GPIO\_OUT1\_DATA 中的相应位会清零。（只写） 
\end{reglist}\end{regdesc}
\end{register}


\begin{register}{H}{GPIO\_ENABLE\_REG}{0x0020}\label{regdesc:GPIOENABLEREG}
\regfield{}{32}{0}{xxxxxxxxxxxxxxxxxxxxxxxxxxxxxxxx}%
\reglabel{Reset}\regnewline%
\begin{regdesc}\begin{reglist}
\item [GPIO\_ENABLE\_REG]GPIO0-31 输出使能。（读/写）
\end{reglist}\end{regdesc}
\end{register}


\begin{register}{H}{GPIO\_ENABLE\_W1TS\_REG}{0x0024}\label{regdesc:GPIOENABLEW1TSREG}
\regfield{}{32}{0}{xxxxxxxxxxxxxxxxxxxxxxxxxxxxxxxx}%
\reglabel{Reset}\regnewline%
\begin{regdesc}\begin{reglist}
\item [GPIO\_ENABLE\_W1TS\_REG]GPIO0-31 输出使能置位寄存器。每一位置 1，则 GPIO\_ENABLE 中的相应位也置 1。（只写）
\end{reglist}\end{regdesc}
\end{register}


\begin{register}{H}{GPIO\_ENABLE\_W1TC\_REG}{0x0028}\label{regdesc:GPIOENABLEW1TCREG}
\regfield{}{32}{0}{xxxxxxxxxxxxxxxxxxxxxxxxxxxxxxxx}%
\reglabel{Reset}\regnewline%
\begin{regdesc}\begin{reglist}
\item [GPIO\_ENABLE\_W1TC\_REG]GPIO0-31 输出使能清零寄存器。每一位置 1，则 GPIO\_ENABLE 中的相应位会清零。（只写） 
\end{reglist}\end{regdesc}
\end{register}


\begin{register}{H}{GPIO\_ENABLE1\_REG}{0x002c}\label{regdesc:GPIOENABLE1REG}
\regfield{(reserved)}{24}{8}{000000000000000000000000}%
\regfield{GPIO\_ENABLE\_DATA}{8}{0}{xxxxxxxx}%
\reglabel{Reset}\regnewline%
\begin{regdesc}\begin{reglist}
\item [GPIO\_ENABLE\_DATA]GPIO32-39 输出使能。（读/写） 
\end{reglist}\end{regdesc}
\end{register}


\begin{register}{H}{GPIO\_ENABLE1\_W1TS\_REG}{0x0030}\label{regdesc:GPIOENABLE1W1TSREG}
\regfield{(reserved)}{24}{8}{000000000000000000000000}%
\regfield{GPIO\_ENABLE\_DATA}{8}{0}{xxxxxxxx}%
\reglabel{Reset}\regnewline%
\begin{regdesc}\begin{reglist}
\item [GPIO\_ENABLE\_DATA]GPIO32-39 输出使能置位寄存器。每一位置 1，则 GPIO\_ENABLE1 中的相应位也置 1。（只写）
\end{reglist}\end{regdesc}
\end{register}


\begin{register}{H}{GPIO\_ENABLE1\_W1TC\_REG}{0x0034}\label{regdesc:GPIOENABLE1W1TCREG}
\regfield{(reserved)}{24}{8}{000000000000000000000000}%
\regfield{GPIO\_ENABLE\_DATA}{8}{0}{xxxxxxxx}%
\reglabel{Reset}\regnewline%
\begin{regdesc}\begin{reglist}
\item [GPIO\_ENABLE\_DATA]GPIO32-39 输出使能清零寄存器。每一位置 1，则 GPIO\_ENABLE1 中的相应位会清零。（只写）
\end{reglist}\end{regdesc}
\end{register}


\begin{register}{H}{GPIO\_STRAP\_REG}{0x0038}\label{regdesc:GPIOSTRAPREG}
\regfield{(reserved)}{16}{16}{0000000000000000}%
\regfield{GPIO\_STRAPPING}{16}{0}{xxxxxxxxxxxxxxxx}%
\reglabel{Reset}\regnewline%
\begin{regdesc}\begin{reglist}
\item [GPIO\_STRAPPING]GPIO strapping 结果：boot\_sel\_chip[5:0] 的 bit5 到 bit0 分别对应 MTDI, GPIO0, GPIO2, GPIO4, MTDO, GPIO5。
\end{reglist}\end{regdesc}
\end{register}


\begin{register}{H}{GPIO\_IN\_REG}{0x003c}\label{regdesc:GPIOINREG}
\regfield{}{32}{0}{xxxxxxxxxxxxxxxxxxxxxxxxxxxxxxxx}%
\reglabel{Reset}\regnewline%
\begin{regdesc}\begin{reglist}
\item [GPIO\_IN\_REG]GPIO0-31 输入值。每个 bit 代表管脚的片外输入值，比如片外引脚为高电平，此 bit 值应为 1，片外引脚为低电平，此 bit 值应为 0。（只读）
\end{reglist}\end{regdesc}
\end{register}


\begin{register}{H}{GPIO\_IN1\_REG}{0x0040}\label{regdesc:GPIOIN1REG}
\regfield{(reserved)}{24}{8}{000000000000000000000000}%
\regfield{GPIO\_IN\_DATA\_NEXT}{8}{0}{xxxxxxxx}%
\reglabel{Reset}\regnewline%
\begin{regdesc}\begin{reglist}
\item [GPIO\_IN\_DATA\_NEXT]GPIO32-39 输入值。每个 bit 代表管脚的片外输入值。（只读）
\end{reglist}\end{regdesc}
\end{register}


\begin{register}{H}{GPIO\_STATUS\_REG}{0x0044}\label{regdesc:GPIOSTATUSREG}
\regfield{GPIO\_STATUS\_INT}{32}{0}{xxxxxxxxxxxxxxxxxxxxxxxxxxxxxxxx}%
\reglabel{Reset}\regnewline%
\begin{regdesc}\begin{reglist}
\item [GPIO\_STATUS\_INT]GPIO0-31 中断状态寄存器。每个 bit 都可以作为两个 CPU 的两种中断源，同时应该把 \hyperref[regdesc:GPIOPINnREG]{GPIO\_PIN\textcolor{red}{\it{n}}\_REG} 的 13-16 bit 相应的 GPIO\_PIN\textcolor{red}{\it{n}}\_INT\_ENA 的使能位置为 1。（读/写）
\end{reglist}\end{regdesc}
\end{register}


\begin{register}{H}{GPIO\_STATUS\_W1TS\_REG}{0x0048}\label{regdesc:GPIOSTATUSW1TSREG}
\regfield{GPIO\_STATUS\_INT\_W1TS}{32}{0}{xxxxxxxxxxxxxxxxxxxxxxxxxxxxxxxx}%
\reglabel{Reset}\regnewline%
\begin{regdesc}\begin{reglist}
\item [GPIO\_STATUS\_INT\_W1TS]GPIO0-31 中断状态置位寄存器。每一位置 1，则 GPIO\_STATUS\_INT 中的相应位也置 1。（只写）
\end{reglist}\end{regdesc}
\end{register}


\begin{register}{H}{GPIO\_STATUS\_W1TC\_REG}{0x004c}\label{regdesc:GPIOSTATUSW1TCREG}
\regfield{GPIO\_STATUS\_INT\_W1TC}{32}{0}{xxxxxxxxxxxxxxxxxxxxxxxxxxxxxxxx}%
\reglabel{Reset}\regnewline%
\begin{regdesc}\begin{reglist}
\item [GPIO\_STATUS\_INT\_W1TC]GPIO0-31 中断状态清除寄存器。每一位置 1，则 GPIO\_STATUS\_INT 中的相应位会清零。（只写）
\end{reglist}\end{regdesc}
\end{register}


\begin{register}{H}{GPIO\_STATUS1\_REG}{0x0050}\label{regdesc:GPIOSTATUS1REG}
\regfield{(reserved)}{24}{8}{000000000000000000000000}%
\regfield{GPIO\_STATUS1\_INT}{8}{0}{xxxxxxxx}%
\reglabel{Reset}\regnewline%
\begin{regdesc}\begin{reglist}
\item [GPIO\_STATUS1\_INT]GPIO32-39 中断状态寄存器。每个 bit 都可以作为两个 CPU 的两种中断源，同时应该把 \hyperref[regdesc:GPIOPINnREG]{GPIO\_PIN\textcolor{red}{\it{n}}\_REG} 的 13-16 bit 相应的 GPIO\_PIN\textcolor{red}{\it{n}}\_INT\_ENA 的使能位置为 1。（读/写）
\end{reglist}\end{regdesc}
\end{register}


\begin{register}{H}{GPIO\_STATUS1\_W1TS\_REG}{0x0054}\label{regdesc:GPIOSTATUS1W1TSREG}
\regfield{(reserved)}{24}{8}{000000000000000000000000}%
\regfield{GPIO\_STATUS1\_INT\_W1TS}{8}{0}{xxxxxxxx}%
\reglabel{Reset}\regnewline%
\begin{regdesc}\begin{reglist}
\item [GPIO\_STATUS1\_INT\_W1TS]GPIO32-39 中断状态置位寄存器。每一位置 1，则 GPIO\_STATUS1\_INT 中的相应位也置 1。（只写）
\end{reglist}\end{regdesc}
\end{register}


\begin{register}{H}{GPIO\_STATUS1\_W1TC\_REG}{0x0058}\label{regdesc:GPIOSTATUS1W1TCREG}
\regfield{(reserved)}{24}{8}{000000000000000000000000}%
\regfield{GPIO\_STATUS1\_INT\_W1TC}{8}{0}{xxxxxxxx}%
\reglabel{Reset}\regnewline%
\begin{regdesc}\begin{reglist}
\item [GPIO\_STATUS1\_INT\_W1TC]GPIO32-39 中断状态清除寄存器。每一位置 1，则 GPIO\_STATUS1\_INT 中的相应位会清零。（只写）
\end{reglist}\end{regdesc}
\end{register}


\begin{register}{H}{GPIO\_ACPU\_INT\_REG}{0x0060}\label{regdesc:GPIOACPUINTREG}
\regfield{}{32}{0}{xxxxxxxxxxxxxxxxxxxxxxxxxxxxxxxx}%
\reglabel{Reset}\regnewline%
\begin{regdesc}\begin{reglist}
\item [GPIO\_ACPU\_INT\_REG]GPIO0-31 APP CPU 中断状态寄存器。（只读）
\end{reglist}\end{regdesc}
\end{register}


\begin{register}{H}{GPIO\_ACPU\_NMI\_INT\_REG}{0x0064}\label{regdesc:GPIOACPUNMIINTREG}
\regfield{}{32}{0}{xxxxxxxxxxxxxxxxxxxxxxxxxxxxxxxx}%
\reglabel{Reset}\regnewline%
\begin{regdesc}\begin{reglist}
\item [GPIO\_ACPU\_NMI\_INT\_REG]GPIO0-31 APP CPU 非屏蔽中断状态寄存器。（只读）
\end{reglist}\end{regdesc}
\end{register}


\begin{register}{H}{GPIO\_PCPU\_INT\_REG}{0x0068}\label{regdesc:GPIOPCPUINTREG}
\regfield{}{32}{0}{xxxxxxxxxxxxxxxxxxxxxxxxxxxxxxxx}%
\reglabel{Reset}\regnewline%
\begin{regdesc}\begin{reglist}
\item [GPIO\_PCPU\_INT\_REG]GPIO0-31 PRO CPU 中断状态寄存器。（只读）
\end{reglist}\end{regdesc}
\end{register}


\begin{register}{H}{GPIO\_PCPU\_NMI\_INT\_REG}{0x006c}\label{regdesc:GPIOPCPUNMIINTREG}
\regfield{}{32}{0}{xxxxxxxxxxxxxxxxxxxxxxxxxxxxxxxx}%
\reglabel{Reset}\regnewline%
\begin{regdesc}\begin{reglist}
\item [GPIO\_PCPU\_NMI\_INT\_REG]GPIO0-31 PRO CPU 非屏蔽中断状态寄存器。（只读）
\end{reglist}\end{regdesc}
\end{register}


\begin{register}{H}{GPIO\_ACPU\_INT1\_REG}{0x0074}\label{regdesc:GPIOACPUINT1REG}
\regfield{(reserved)}{24}{8}{000000000000000000000000}%
\regfield{GPIO\_APPCPU\_INT}{8}{0}{xxxxxxxx}%
\reglabel{Reset}\regnewline%
\begin{regdesc}\begin{reglist}
\item [GPIO\_APPCPU\_INT]GPIO32-39 APP CPU 中断状态寄存器。（只读）
\end{reglist}\end{regdesc}
\end{register}


\begin{register}{H}{GPIO\_ACPU\_NMI\_INT1\_REG}{0x0078}\label{regdesc:GPIOACPUNMIINT1REG}
\regfield{(reserved)}{24}{8}{000000000000000000000000}%
\regfield{GPIO\_APPCPU\_NMI\_INT}{8}{0}{xxxxxxxx}%
\reglabel{Reset}\regnewline%
\begin{regdesc}\begin{reglist}
\item [GPIO\_APPCPU\_NMI\_INT]GPIO32-39 APP CPU 非屏蔽中断状态寄存器。（只读）
\end{reglist}\end{regdesc}
\end{register}


\begin{register}{H}{GPIO\_PCPU\_INT1\_REG}{0x007c}\label{regdesc:GPIOPCPUINT1REG}
\regfield{(reserved)}{24}{8}{000000000000000000000000}%
\regfield{GPIO\_PROCPU\_INT}{8}{0}{xxxxxxxx}%
\reglabel{Reset}\regnewline%
\begin{regdesc}\begin{reglist}
\item [GPIO\_PROCPU\_INT]GPIO32-39 PRO CPU 中断状态寄存器。（只读）
\end{reglist}\end{regdesc}
\end{register}


\begin{register}{H}{GPIO\_PCPU\_NMI\_INT1\_REG}{0x0080}\label{regdesc:GPIOPCPUNMIINT1REG}
\regfield{(reserved)}{24}{8}{000000000000000000000000}%
\regfield{GPIO\_PROCPU\_NMI\_INT}{8}{0}{xxxxxxxx}%
\reglabel{Reset}\regnewline%
\begin{regdesc}\begin{reglist}
\item [GPIO\_PROCPU\_NMI\_INT]GPIO32-39 PRO CPU 非屏蔽中断状态寄存器。（只读）
\end{reglist}\end{regdesc}
\end{register}


\begin{register}{H}{GPIO\_PIN\textcolor{red}{\it{n}}\_REG (\textcolor{red}{\it{n}}: 0-19, 21-23, 25-27, 32-39)}{0x88+0x4*\textcolor{red}{\it{n}}}\label{regdesc:GPIOPINnREG}
\regfield{(reserved)}{14}{18}{00000000000000}%
\regfield{GPIO\_PIN\textcolor{red}{\it{n}}\_INT\_ENA}{5}{13}{xxxxx}%
\regfield{(reserved)}{2}{11}{00}%
\regfield{GPIO\_PIN\textcolor{red}{\it{n}}\_WAKEUP\_ENABLE}{1}{10}{x}%
\regfield{GPIO\_PIN\textcolor{red}{\it{n}}\_INT\_TYPE}{3}{7}{xxx}%
\regfield{(reserved)}{4}{3}{0000}%
\regfield{GPIO\_PIN\textcolor{red}{\it{n}}\_PAD\_DRIVER}{1}{2}{x}%
\regfield{(reserved)}{2}{0}{00}%
\reglabel{Reset}\regnewline%
\begin{regdesc}\begin{reglist}
\item [GPIO\_PIN\textcolor{red}{\it{n}}\_INT\_ENA]bit0：APP CPU 中断使能；（读/写） \\bit1：APP CPU 非屏蔽中断使能； \\bit2：PRO CPU 中断使能；  \\bit3：PRO CPU 非屏蔽中断使能。 
\item [GPIO\_PIN\textcolor{red}{\it{n}}\_WAKEUP\_ENABLE]GPIO 唤醒使能。只能将 CPU 从 Light-sleep 中唤醒。（读/写）
\item [GPIO\_PIN\textcolor{red}{\it{n}}\_INT\_TYPE]0: GPIO 中断类型；（读/写） \\1：上升沿触发；\\2：下降沿触发；\\3：任一沿触发；\\4：低电平触发；\\5：高电平触发。
\item [GPIO\_PIN\textcolor{red}{\it{n}}\_PAD\_DRIVER]0：正常输出； 1：开漏方式输出。（读/写）
\end{reglist}\end{regdesc}
\end{register}


\begin{register}{H}{GPIO\_FUNC\textcolor{red}{\it{y}}\_IN\_SEL\_CFG\_REG (\textcolor{red}{\it{y}}: 0-255)}{0x130+0x4*\textcolor{red}{\it{y}}}\label{regdesc:GPIOFUNCnINSELCFGREG}
\regfield{(reserved)}{24}{8}{000000000000000000000000}%
\regfield{GPIO\_SIG\textcolor{red}{\it{y}}\_IN\_SEL}{1}{7}{x}%
\regfield{GPIO\_FUNC\textcolor{red}{\it{y}}\_IN\_INV\_SEL}{1}{6}{x}%
\regfield{GPIO\_FUNC\textcolor{red}{\it{y}}\_IN\_SEL}{6}{0}{xxxxxx}%
\reglabel{Reset}\regnewline%
\begin{regdesc}\begin{reglist}
\item [GPIO\_SIG\textcolor{red}{\it{y}}\_IN\_SEL]旁路 GPIO 交换矩阵。1：通过 GPIO 交换矩阵；0：直接通过 IO MUX 连接信号与外设。（读/写）
\item [GPIO\_FUNC\textcolor{red}{\it{y}}\_IN\_INV\_SEL]反转输入值。1:反转；0:不反转。（读/写）
\item [GPIO\_FUNC\textcolor{red}{\it{y}}\_IN\_SEL]外设输入 \textcolor{red}{\it{y}} 选择控制。此位选择 1 个 GPIO 交换矩阵输入管脚中与信号连接，或者选择 0x38 与恒高电平输入信号连接或者选择 0x30 与恒低电平输入信号连接。（读/写） 
\end{reglist}\end{regdesc}
\end{register}


\begin{register}{H}{GPIO\_FUNC\textcolor{red}{\it{n}}\_OUT\_SEL\_CFG\_REG (\textcolor{red}{\it{n}}: 0-19, 21-23, 25-27, 32-33)}{0x530+0x4*\textcolor{red}{\it{n}}}\label{regdesc:GPIOFUNCnOUTSELCFGREG}
\regfield{(reserved)}{20}{12}{00000000000000000000}%
\regfield{GPIO\_FUNC\textcolor{red}{\it{n}}\_OEN\_INV\_SEL}{1}{11}{x}%
\regfield{GPIO\_FUNC\textcolor{red}{\it{n}}\_OEN\_SEL}{1}{10}{x}%
\regfield{GPIO\_FUNC\textcolor{red}{\it{n}}\_OUT\_INV\_SEL}{1}{9}{x}%
\regfield{GPIO\_FUNC\textcolor{red}{\it{n}}\_OUT\_SEL}{9}{0}{xxxxxxxxx}%
\reglabel{Reset}\regnewline%
\begin{regdesc}\begin{reglist}
\item [GPIO\_FUNC\textcolor{red}{\it{n}}\_OEN\_INV\_SEL]1：反转输出使能信号；0：不反转输出使能信号。（读/写）
\item \label{fielddesc:GPIOFUNCNOENSEL} [GPIO\_FUNC\textcolor{red}{\it{n}}\_OEN\_SEL]1：强制从 GPIO\_ENABLE\_REG bit \textcolor{red}{\it{n}} 中选取输出使能信号。0：使用从外设选取的输出使能信号。（读/写）
\item [GPIO\_FUNC\textcolor{red}{\it{n}}\_OUT\_INV\_SEL]1：反转输出值；0：不反转输出值。（读/写）
\item [GPIO\_FUNC\textcolor{red}{\it{n}}\_OUT\_SEL] GPIO 输出 \textcolor{red}{\it{n}} 的选择控制。值为  \textcolor{red}{\it{s}} (0<=\textcolor{red}{\it{s}}<256) 连接外设输出 \textcolor{red}{\it{s}} 与 GPIO 输出 \textcolor{red}{\it{n}}。值为 256 选择 GPIO\_OUT\_REG/GPIO\_OUT1\_REG 和 GPIO\_ENABLE\_REG/GPIO\_ENABLE1\_REG 的 bit \textcolor{red}{\it{n}} 做为输出值和输出使能。（读/写）
\end{reglist}\end{regdesc}
\end{register}
